% =====================================================================
%  2bat-ud6-6.tex  —  SESSIÓ 6: Interpretar intervals de confiança (mecanització)
%  (sense preàmbul: s'inclou des de main.tex)
% =====================================================================
\horatitol{Sessió 6 — Interpretar intervals de confiança}%
{Mecanitzar la construcció i la interpretació d'intervals de confiança amb una bateria graduada i correcció creuada.}

\subsection{Idea clau}

Ja sabeu què és un interval de confiança. Avui el \textbf{mecanitzem}: construir-lo a
partir d'un resultat amb marge, expressar amb precisió què afirma (i què no), i esquivar
els errors d'interpretació. És una de les habilitats que més s'avaluen i la més útil a la
vida quotidiana.

\begin{idea}
\textbf{Construir i interpretar l'interval, pas a pas}
\begin{enumerate}[leftmargin=1.6em, itemsep=2pt]
  \item \textbf{Construir:} interval $=$ resultat $\pm$ marge (resta i suma el marge).
  \item \textbf{Interpretar:} «és molt probable que el valor real de \emph{tota la
        població} estigui dins d'aquest interval».
  \item \textbf{Nivell de confiança} (sovint $95\%$): indica \emph{com de segur} és el
        mètode, \textbf{no} que «el 95\% de la gent» opini res.
\end{enumerate}
\textbf{Errors a evitar:} (a) tractar l'interval com una certesa; (b) dir «el 95\% de la
gent» quan el $95\%$ és el nivell de confiança; (c) comparar dos resultats sense mirar si
els intervals \textbf{se superposen}.
\end{idea}

\subsection{Exemples}

\begin{exemple}[construir i dir-ho bé]
Resultat: «$63\%\pm 3\%$, amb un nivell de confiança del $95\%$». Interval:
$[60\%,\,66\%]$. \textbf{Lectura correcta:} «amb un $95\%$ de confiança, la proporció real
a tota la població està entre el $60\%$ i el $66\%$». \textbf{Lectura incorrecta:} «el
$95\%$ de la població opina X» (confon el nivell de confiança amb un percentatge de gent).
\end{exemple}

\begin{exemple}[comparar dos resultats amb els seus intervals]
Un estudi diu que el servei A té un $52\%\pm 4\%$ de satisfacció ($[48\%,56\%]$) i el B un
$55\%\pm 4\%$ ($[51\%,59\%]$). Els intervals \textbf{se superposen} (del $51\%$ al $56\%$),
així que \textbf{no} es pot afirmar que B sigui millor que A: la diferència podria ser fruit
de l'atzar del mostreig.
\end{exemple}

\subsection{Exercicis resolts}

\begin{exresolt}[interval i interpretació correcta]
Una enquesta dona «$47\%\pm 3\%$» amb confiança del $95\%$. Construeix l'interval i escriu
la interpretació correcta.
\solucio
Interval: $47\%-3\%=44\%$ fins a $47\%+3\%=50\%$, és a dir $[44\%,\,50\%]$.
\textbf{Interpretació:} amb un $95\%$ de confiança, la proporció real a tota la població
està entre el $44\%$ i el $50\%$. No es pot afirmar que sigui exactament el $47\%$.
\resposta{$[44\%,\,50\%]$; el valor real hi és molt probablement (confiança $95\%$).}
\end{exresolt}

\begin{exresolt}[superposició d'intervals]
Dues edicions d'un programa donen un èxit del $60\%\pm 5\%$ i del $64\%\pm 5\%$. Es pot dir
que la segona edició ha estat millor?
\solucio
Intervals: $[55\%,65\%]$ i $[59\%,69\%]$. \textbf{Se superposen} (del $59\%$ al $65\%$),
així que \textbf{no} es pot afirmar amb seguretat que la segona sigui millor: la diferència
de $4$ punts cau dins del que podria explicar l'atzar del mostreig.
\resposta{No; els intervals se superposen, la diferència pot ser casual.}
\end{exresolt}

\subsection{Exercicis per fer}

\noindent\textit{Resoleu i, en acabar, intercanvieu el full amb una altra parella per a la
correcció creuada.}

\begin{exercicis}
\begin{exlist}
  \item Construeix l'interval de «$41\%\pm 3\%$» i escriu què afirma correctament l'estudi.
  \item Construeix l'interval de «$28\%\pm 4\%$».
  \item Digues si és correcta la frase: «amb un nivell de confiança del $95\%$, el $95\%$
        de la població opina X». Si no ho és, corregeix-la.
  \item Dos serveis tenen una satisfacció del $70\%\pm 3\%$ i del $72\%\pm 3\%$. Se
        superposen els intervals? Es pot dir que un és millor que l'altre?
  \item \textbf{(Comparació)} Un programa passa d'un $45\%\pm 2\%$ a un $53\%\pm 2\%$
        d'èxit. Se superposen els intervals? Es pot afirmar que ha millorat?
  \item \textbf{(Interpretació)} Explica la diferència entre el \emph{marge d'error} i el
        \emph{nivell de confiança} d'una enquesta.
\end{exlist}
\end{exercicis}

% ---------------------------------------------------------------------
\fullsolucions{Sessió 6}
\begin{solucions}
\begin{sollist}
  \item $[38\%,\,44\%]$. És molt probable que el valor real de la població estigui entre el
        $38\%$ i el $44\%$ (no exactament el $41\%$).
  \item $28\%-4\%=24\%$ fins a $28\%+4\%=32\%$: interval $[24\%,\,32\%]$.
  \item Incorrecta: confon el nivell de confiança amb un percentatge de gent. Correcte:
        «amb un $95\%$ de confiança, el valor real de la població està dins de l'interval»;
        el $95\%$ mesura la fiabilitat del mètode, no quanta gent opina X.
  \item Intervals $[67\%,73\%]$ i $[69\%,75\%]$: \textbf{se superposen} (del $69\%$ al
        $73\%$). No es pot afirmar que un sigui millor; la diferència pot ser casual.
  \item Intervals $[43\%,47\%]$ i $[51\%,55\%]$: \textbf{no} se superposen. Sí que es pot
        afirmar que ha millorat: la diferència difícilment s'explica per l'atzar.
  \item Resposta oberta. Idea: el \textbf{marge d'error} diu com d'ampli és l'interval al
        voltant del resultat (la precisió de l'estimació); el \textbf{nivell de confiança}
        diu amb quina seguretat el valor real cau dins d'aquest interval (habitualment
        $95\%$). Són coses diferents.
\end{sollist}
\end{solucions}
